\section{Síntesis y ampliación}

Esta clase puede condensarse en una sola frase: \textbf{entrenar un helpful chatbot no consiste en elegir un optimizer, sino en diseñar de forma conjunta la data distribution, el preference protocol, el training objective y el measurement system.} Los H4 human-data experiments y el Zephyr synthetic-data pipeline muestran dos rutas viables; la GPT-4 evaluator audit demuestra, a su vez, que si la measurement y los training data tienen el mismo origen, el ranking premia el style matching y no el objetivo real.

\subsection{Once conclusiones aplicables}

\begingroup
\footnotesize
\setlength{\columnsep}{1.2em}
\begin{multicols}{2}
\begin{enumerate}[itemsep=0.05em,topsep=0.1em,parsep=0pt,leftmargin=*]
  \item Reportar la dataset provenance: quién escribe el prompt, quién escribe la completion y quién filtra.
  \item Describir los datos con la task/length/turn distribution, no solo con el example count.
  \item Distinguir las unidades de conteo de demonstration, preference pair y policy rollout.
  \item El SFT establece primero el instruction-following; DPO/RLHF suele ser un refinement, no un sustituto de esa base.
  \item Los human-written data y los synthetic data tienen errores con correlaciones distintas; conviene conservar un holdout heterogéneo.
  \item El reward model debe evaluarse con nuevos generators, pairs difíciles de distinguir y held-out tasks.
  \item Reportar a la vez task metrics, human preference, LLM judge y response statistics.
  \item Ante una metric reversal, no se elige un ganador: se localiza la evaluator preference.
  \item El LLM judge requiere order swap, length control, human calibration y task stratification.
  \item El red teaming mide los riesgos de cola; no puede sustituirse por un helpfulness score promedio.
  \item El costo, la escala y el leaderboard deben indicar siempre fecha, unidades y protocol.
\end{enumerate}
\end{multicols}
\endgroup

\begin{importantbox}{Visión final del sistema}
Un alignment pipeline puede escribirse como un ciclo cerrado:
\[
\text{Data source}\rightarrow\text{Training objective}\rightarrow\text{Candidate behavior}
\rightarrow\text{Evaluator}\rightarrow\text{Next data round}.
\]
El sesgo del Evaluator determina qué datos se recolectan en la siguiente ronda, por lo que el measurement error vuelve a escribirse en la training distribution. El verdadero objetivo de ingeniería no es ganar el ranking una vez, sino lograr que el ciclo cerrado pueda auditarse, calibrarse y corregirse.
\end{importantbox}

\subsection{Lecturas adicionales}

\begingroup
\scriptsize
\begin{itemize}[itemsep=0pt,topsep=0.15em,parsep=0pt]
  \item Ouyang et al., \href{https://arxiv.org/abs/2203.02155}{Training language models to follow instructions with human feedback}: flujo básico de SFT, reward modeling y PPO/RLHF.
  \item Wang et al., \href{https://arxiv.org/abs/2212.10560}{Self-Instruct}; Ding et al., \href{https://arxiv.org/abs/2305.14233}{UltraChat}; Li et al., \href{https://arxiv.org/abs/2303.17760}{CAMEL}: tres rutas de synthetic instruction/dialogue generation.
  \item Tunstall et al., \href{https://arxiv.org/abs/2310.16944}{Zephyr: Direct Distillation of LM Alignment}: dSFT, AI feedback y dDPO.
  \item Zheng et al., \href{https://arxiv.org/abs/2306.05685}{Judging LLM-as-a-Judge with MT-Bench and Chatbot Arena}: LLM judge, position bias y human agreement.
  \item Hugging Face H4, \href{https://huggingface.co/blog/llm-leaderboard}{LLM evaluator bias analysis} y \href{https://huggingface.co/blog/red-teaming}{red-teaming workflow}: punto de acceso oficial a la práctica detrás de las slides de la clase.
\end{itemize}
\endgroup

\end{document}
